\section*{अध्याय \ref{ch:g12:domesticated-plants} --- पालतू बनाया गया पौधा}

\begin{solution}{exo:g12:domesticated-plants:1}
\omterm{def:g12:domesticated-plants:domestication}{पालतूकरण}: किसी जंगली \omterm{def:g10:biodiversity-scales:species}{जाति} का ऐसी \omterm{def:g10:biodiversity-scales:species}{जाति} में बदलना जो मनुष्यों पर निर्भर हो और
उनकी सेवा करे। \omterm{def:g12:domesticated-plants:domestication}{कृत्रिम वरण}: मनुष्यों का उन्हीं व्यक्तियों को रखना और बोना
जिनमें उनके चाहे लक्षण हों, जिससे उनके पीछे बैठे \omterm{def:g10:universal-dna:gene}{युग्मविकल्पियों} की
बारंबारता चढ़ती है।
\end{solution}

\begin{solution}{exo:g12:domesticated-plants:2}
न झड़ती बालियाँ (कोई प्रकीर्णन नहीं); तुरंत अंकुरण (किसी बुरे साल को टालने
का कोई ज़रिया नहीं); विषों का खो जाना (हर कोई खा जाता है); बिना शाखा वाले
पौधे पर एक बड़ी बाली (सारे अंडे एक ही टोकरी में); और एक साथ पकना (एक बुरा
हफ़्ता पूरी फ़सल तबाह कर देता है)।
\end{solution}

\begin{solution}{exo:g12:domesticated-plants:3}
दो समयुग्मजी अंतःप्रजनित वंशक्रमों की एकसार, ओजस्वी संतान। उसका अपना बीज कम
ओजस्वी पौधों के मिश्रण में छँट जाता है, इसलिए वह किसान हर साल ताज़ा संकर
बीज ख़रीदता है।
\end{solution}

\begin{solution}{exo:g12:domesticated-plants:4}
वह \omterm{def:g10:universal-dna:gene}{जीन} उस \omterm{def:g10:universal-dna:information}{DNA} में डाला जाता है जो मिट्टी का कोई जीवाणु प्राकृतिक रूप से पादप
\omterm{def:g10:cells-common-unit:cell}{कोशिकाओं} में पहुँचाता है; और जिन संक्रमित \omterm{def:g10:cells-common-unit:cell}{कोशिकाओं} ने वह \omterm{def:g10:universal-dna:gene}{जीन} लिया उनसे पूरे
पौधे दोबारा उगाए जाते हैं।
\end{solution}

\begin{solution}{exo:g12:domesticated-plants:5}
लगभग \num{1.9} से \qty{7.2}{t/ha} तक: यानी लगभग चार गुना।
\end{solution}

\begin{solution}{exo:g12:domesticated-plants:6}
जिन पौधों का दाना झड़ जाता था वे उसे कटाई से पहले या उसी के दौरान खो देते
थे; और सख़्त डंठल वाले उत्परिवर्ती अपना दाना थामे रहते थे तथा वही बटोरे
गए, यानी वही बोए गए। कटाई ने ही, पीढ़ी दर पीढ़ी, सख़्त डंठल वाला
\omterm{def:g10:universal-dna:gene}{युग्मविकल्पी} चुना।
\end{solution}

\begin{solution}{exo:g12:domesticated-plants:7}
एक ही \omterm{def:g10:biodiversity-scales:species}{जाति}: वे आपस में प्रजनन करती हैं और उर्वर संतान देती हैं। किसानों ने
हर वंशक्रम में किसी अलग अंग के बढ़ावे को चुना --- पत्ते, शीर्ष कली, फूल की
कलियाँ --- इसलिए वे बढ़त के नियमन में अलग हैं, आपस में संकरित होने की अपनी
क्षमता में नहीं।
\end{solution}

\begin{solution}{exo:g12:domesticated-plants:8}
F1 सब $Aa$। F2: $AA$ $1/4$, $Aa$ $1/2$, $aa$ $1/4$। F1 इसलिए एकसार है
क्योंकि हर पौधे को एक जनक से $A$ मिलता है और दूसरे से $a$; जबकि F2 में
\omterm{def:g12:meiosis-genetic-shuffling:meiosis}{अर्धसूत्री विभाजन} और निषेचन \omterm{def:g10:universal-dna:gene}{युग्मविकल्पियों} को यादृच्छिक ढंग से बाँट देते
हैं।
\end{solution}

\begin{solution}{exo:g12:domesticated-plants:9}
लाखों छेदकों में से कुछ ऐसा \omterm{def:g11:mutations:mutation}{उत्परिवर्तन} ढोते थे जो उन्हें वह विष झेलने देता
था; और जिस खेत में हर पौधा उसे बनाता हो वहाँ केवल वही बचे तथा प्रजनन कर पाए
--- यानी वरण। कोई शरणस्थल बहुत-से संवेदनशील छेदक ज़िंदा रखता है; और उनके
दुर्लभ प्रतिरोधियों से हुए संगम ऐसी संतान देते हैं जो एक ही प्रतिरोध
\omterm{def:g10:universal-dna:gene}{युग्मविकल्पी} ढोती है, जो (\omterm{def:g11:genes-and-disease:genetic}{अप्रभावी} होने के कारण) विष वाली मक्का पर मर जाती
है, इसलिए वह \omterm{def:g10:universal-dna:gene}{युग्मविकल्पी} हर पीढ़ी में पतला होता जाता है।
\end{solution}

\begin{solution}{exo:g12:domesticated-plants:10}
इनसे जुड़े \omterm{def:g10:universal-dna:gene}{जीन} परिवर्धन को नियंत्रित करते हैं --- पौधे का शाखन, दाने के कवच
की बढ़त, क़तारों की संख्या; और ऐसे किसी \omterm{def:g10:universal-dna:gene}{जीन} के नियमन में आया कोई बदलाव पूरे
अंग का रूप बदल देता है। कम नियामक बदलाव, बड़े दिखते असर।
\end{solution}

\begin{solution}{exo:g12:domesticated-plants:11}
वे उन \omterm{def:g10:universal-dna:gene}{युग्मविकल्पियों} को थामे हैं जिन्हें वरण ने उस फ़सल से हटा दिया:
बीमारियों तथा कीटों के प्रति प्रतिरोध, सूखे, नमक या ठंड की सहनशीलता। कोई नई
बीमारी या जलवायु आए तो प्रजनक उन \omterm{def:g10:universal-dna:gene}{युग्मविकल्पियों} को संकरण से उस फ़सल में
वापस लाते हैं; और कोई जंगली नातेदार एक बार विलुप्त हो जाए तो उसके
\omterm{def:g10:universal-dna:gene}{युग्मविकल्पी} जा चुके होते हैं।
\end{solution}

\begin{solution}{exo:g12:domesticated-plants:12}
उस तोरिया का पराग उस सरसों को निषेचित करता है; संकर वह सहनशीलता \omterm{def:g10:universal-dna:gene}{जीन} ढोते
हैं; कुछ पीढ़ियों पर सरसों से प्रति-संकरण और खेतों पर छिड़के शाकनाशी से
वरण होने पर कोई सहनशील जंगली सरसों फैल जाती है। प्रक्रियाएँ: संकरण तथा
\omterm{def:g12:selection-drift-speciation:population}{आबादियों} के बीच किसी \omterm{def:g10:universal-dna:gene}{जीन} की क्षैतिज हलचल, और फिर उस शाकनाशी से \omterm{def:g12:selection-drift-speciation:selection}{प्राकृतिक
वरण}।
\end{solution}

\begin{solution}{exo:g12:domesticated-plants:13}
जैविक: उस पौधे का \omterm{def:g10:cell-metabolism:autotrophy}{प्रकाशसंश्लेषण} और दाना भरने की उसकी क्षमता अपनी सीमाओं के
पास हैं, और उन उत्तम क़िस्मों की \omterm{def:g10:biodiversity-scales:biodiversity}{आनुवंशिक विविधता}, जिनसे आगे के लाभ आने
हैं, सँकरी है। ग़ैर-जैविक: उर्वरक तथा कीटनाशी का इस्तेमाल अब बढ़ नहीं रहा,
या घटाया जा रहा है, और जलवायु के दबाव बढ़े हैं। कोई प्रजनक जंगली नातेदारों
से \omterm{def:g10:universal-dna:gene}{युग्मविकल्पी} ला सकता है, या \omterm{def:g10:cell-metabolism:autotrophy}{प्रकाशसंश्लेषण} अथवा सूखा-सहनशीलता के \omterm{def:g10:universal-dna:gene}{जीन}
संपादित कर सकता है।
\end{solution}

\begin{solution}{exo:g12:domesticated-plants:14}
संकरण: वह \omterm{def:g10:universal-dna:gene}{युग्मविकल्पी} किसी ऐसे पौधे में होना चाहिए जिसका उस फ़सल से संकरण
हो सके; हज़ारों \omterm{def:g10:universal-dna:gene}{जीन} फिर से फेंटे जाते हैं और उन्हें कोई दस पीढ़ियों पर वरण
से छाँटना पड़ता है। \omterm{prop:g12:domesticated-plants:transgenic}{परजीन-प्रवेश}: वह \omterm{def:g10:universal-dna:gene}{जीन} किसी भी जीव से आ सकता है; एक ही \omterm{def:g10:universal-dna:gene}{जीन}
जोड़ा जाता है और बाक़ी जीनोम अछूता रहता है; और कुछ ही साल लगते हैं।
\end{solution}

\begin{solution}{exo:g12:domesticated-plants:15}
अंकगणित बिलकुल वही है: कुछ \omterm{def:g10:universal-dna:gene}{युग्मविकल्पी} ढोते व्यक्ति अधिक प्रजनन करते हैं,
और उन \omterm{def:g10:universal-dna:gene}{युग्मविकल्पियों} की बारंबारताएँ चढ़ती हैं। जो अलग है वह है वह कारक ---
पर्यावरण की जगह कोई मानव चुनाव --- और उस पौधे के लिए नतीजा: \omterm{def:g12:selection-drift-speciation:selection}{प्राकृतिक वरण}
किसी \omterm{def:g10:biodiversity-scales:species}{जाति} को अपने बल पर जीने लायक़ बनाए रखता है, जबकि \omterm{def:g12:domesticated-plants:domestication}{पालतूकरण} ऐसे लक्षण
चुनता है जो चुनने वाले के काम आएँ और उस पौधे को उसी पर निर्भर छोड़ देता है।
\omterm{def:g12:domesticated-plants:domestication}{पालतूकरण} किसी दूसरे चुनने वाले के साथ \omterm{def:g12:selection-drift-speciation:selection}{प्राकृतिक वरण} है, उसका उल्टा नहीं।
\end{solution}

\begin{solution}{pb:g12:domesticated-plants:1}
\textbf{1.} टिओसिंटे $12 \times 0.08 \approx \qty{1}{g}$; मक्का $600
\times 0.3 = \qty{180}{g}$: यानी लगभग दो सौ गुना।

\textbf{2.} किसी कड़ी गुल्ली पर बहुत-से दाने (कोई प्रकीर्णन नहीं); नंगे दाने
(जानवरों तथा सड़न से असुरक्षित); एक ही डंठल पर एक बड़ी बाली (जोखिम का कोई
बँटवारा नहीं); छिलकों में थमे दाने (ज़मीन पर गिर ही नहीं सकते)।

\textbf{3.} नंगा \omterm{def:g10:universal-dna:gene}{युग्मविकल्पी} \omterm{def:g11:genes-and-disease:genetic}{अप्रभावी} है: कोई पौधा नंगे दाने तभी दिखाता है
जब वह दो प्रतियाँ ढोए, जो किसी दुर्लभ \omterm{def:g10:universal-dna:gene}{युग्मविकल्पी} के लिए पौधों की एक नन्ही
अल्पसंख्या में ही होता है; और एक प्रति के वाहक साधारण दिखते हैं।

\textbf{4.} $q^2 = 0.0001$: दस हज़ार में एक पौधा।

\textbf{5.} सारा बीज $nn$ पौधों से आता है, जिनके दानों को माँ से $n$ मिला;
और पराग ज़्यादातर साधारण खेत से आया ($N$, 0.99 बारंबारता के साथ)। वह बीज
लगभग सारा $Nn$ है, इसलिए $n$ की बारंबारता उछलकर लगभग 0.5 हो जाती है --- यानी
सख़्त वरण की एक ही पीढ़ी।

\textbf{6.} 9000 पीढ़ियों पर $600/12 = 50$: यानी प्रति पीढ़ी $50^{1/9000}
\approx 1.0004$ का गुणक, यानी 0.04\% का लाभ। नन्हे क़दम, जो हज़ारों
पीढ़ियों पर चक्रवृद्धि हों, एक बड़ा बदलाव बन जाते हैं --- यही वरण का
अंकगणित है।

\textbf{7.} जब बहुत-से \omterm{def:g10:universal-dna:gene}{जीन} हर एक थोड़ा-थोड़ा जोड़ते हों, तो वरण उन सब पर
अनुकूल \omterm{def:g10:universal-dna:gene}{युग्मविकल्पियों} की बारंबारता चढ़ाता है और लंबे समय तक नए जोड़ खोजता
रहता है; जबकि \omterm{def:g10:universal-dna:gene}{एक-जीन} वाला लक्षण उसी वक़्त पलट जाता है जब वह \omterm{def:g11:genes-and-disease:genetic}{अप्रभावी}
\omterm{def:g10:universal-dna:gene}{युग्मविकल्पी} समयुग्मजी बना दिया जाए।

\textbf{8.} एक ही \omterm{def:g10:biodiversity-scales:species}{जाति}: उर्वर संकर। \omterm{def:g12:domesticated-plants:domestication}{पालतूकरण} ने कुछ \omterm{def:g10:universal-dna:gene}{जीनों} पर \omterm{def:g10:universal-dna:gene}{युग्मविकल्पियों}
की बारंबारताएँ और उस पौधे का रूप बदला; उसने कोई प्रजनन रुकावट नहीं बनाई।

\textbf{9.} जो बीमारी या कीट एक पौधे को हरा दे वह उन सबको हरा देता है; और
जवाब वह विविधता है जो पुरानी क़िस्मों तथा जंगली नातेदारों में, बीज बैंकों
और खेतों में, रखी है।

\textbf{10.} वह किसान ही उस पौधे का प्रकीर्णन, उसका वरण करता पर्यावरण और
अगली पीढ़ी तक उसका इकलौता रास्ता है: मक्का हर पीढ़ी में \omterm{def:g12:domesticated-plants:domestication}{कृत्रिम वरण} के नीचे
है और उससे निकल नहीं सकती।

\textbf{11.} बेहतर जनक से $10/5.5 - 1 \approx 82\%$ ऊपर; और F1 से
$1 - 8.2/10 = 18\%$ नीचे।

\textbf{12.} हर वंशक्रम कुछ कमज़ोर \omterm{def:g11:genes-and-disease:genetic}{अप्रभावी} \omterm{def:g10:universal-dna:gene}{युग्मविकल्पियों} के लिए
समयुग्मजी है; और F1 में ऐसा हर \omterm{def:g10:universal-dna:gene}{युग्मविकल्पी} दूसरे वंशक्रम के काम करते
\omterm{def:g10:universal-dna:gene}{युग्मविकल्पी} से ढँक जाता है। F2 में ऐसे हर \omterm{def:g10:universal-dna:gene}{जीन} पर एक चौथाई पौधे उस कमज़ोर
\omterm{def:g10:universal-dna:gene}{युग्मविकल्पी} के लिए समयुग्मजी होते हैं, और औसत गिर जाता है।

\textbf{13.} ख़रीदना: $10 - 0.4 = \qty{9.6}{t/ha}$; बचाना: \qty{8.2}{t/ha}।
ख़रीदना हर साल \qty{1.4}{t/ha} से जीतता है, और इसीलिए वह कंपनी जो उन
अंतःप्रजनित वंशक्रमों की मालिक है हर मौसम संकर बीज बेच सकती है।

\textbf{14.} वह कंपनी उन वंशक्रमों को अपने पास रखती है; उनके साथ कोई किसान
अनिश्चित काल तक वह संकर बना सकता और वह कंपनी एक ही बार बेच पाती। वे वंशक्रम
उस कंपनी की पूँजी हैं; केवल उनका उत्पाद बेचा जाता है।

\textbf{15.} किसी संकर क़िस्म का हर खेत आनुवंशिक रूप से समरूप है, और इलाक़े
कुछ ही संकर उगाते हैं: यानी विविधता सँकरी हो गई। और वह किसान अब बीज नहीं
बनाता तथा हर साल आपूर्तिकर्ता पर निर्भर है।

\textbf{16.} $q^2 = 10^{-6}$: दस लाख में एक छेदक बचता है; और बचे हुओं में,
जो सब $rr$ हैं, $r$ की बारंबारता 1 है।

\textbf{17.} बचे हुए आपस में प्रजनन करते हैं; उनकी संतान सब $rr$ होती है;
और कुछ ही पीढ़ियों में उस खेत के छेदक प्रतिरोधी हो जाते हैं --- ठीक वैसे
जैसे अधूरे कोर्स पर बैठे किसी रोगी में जीवाणु।

\textbf{18.} कोई प्रतिरोधी बचा छेदक ($rr$) किसी शरणस्थल छेदक ($RR$, यानी
भारी बहुमत) से संगम करके $Rr$ संतान देता है, जो विष वाली मक्का पर मर जाती
है। जब तक संवेदनशील छेदक प्रतिरोधियों से कहीं अधिक हों, तब तक लगभग हर
प्रतिरोधी छेदक के बच्चे विषमयुग्मजी होते हैं और हटा दिए जाते हैं, तथा $r$
दुर्लभ बना रहता है।

\textbf{19.} एक जैसा: दोनों का मक़सद उन दुर्लभ प्रतिरोधी व्यक्तियों का गुणन
रोकना है। अलग: पूरा कोर्स उन्हें हर व्यक्ति को मारकर हटाता है, इससे पहले कि
प्रतिरोध चुना जा सके; जबकि शरणस्थल बचे हुओं को क़बूल करता है और उन्हें
संवेदनशील साथियों से पतला कर देता है।

\textbf{20.} प्रति बाली लगभग दो सौ गुना अधिक दाना; लगभग पाँच बड़े \omterm{def:g10:universal-dna:gene}{जीन}; और
ऐसा पौधा जो बिखरने तथा बोए जाने के लिए किसान पर निर्भर है।
\end{solution}
